\documentclass[english, parskip=full]{scrartcl}
\usepackage[utf8]{inputenc}  % Kodierung der Datei
\usepackage[english]{babel}  % Sprache auf Deutsch 
\usepackage{color}
\usepackage{graphicx, gensymb}
\usepackage{amsfonts, cancel, amssymb}
\usepackage{amsmath, amsthm, array, esint}
\usepackage{booktabs, multirow}  % schönere Tabellen
\usepackage{caption}  % Anpassung der Captions
\usepackage{scrlayer-scrpage}    % Manipulation des Seitenstils
\pagestyle{scrheadings}  % Stil für die Seite setzen
\automark{section}  % Abschnittsnamen als Seitenbeschriftung 
\setheadsepline{.5pt}    % Kopzeile durch Linie abtrennen
\usepackage[hidelinks]{hyperref}  % Links und weitere PDF-Features
\usepackage{typearea, tikz, pgffor, verbatim}
\usetikzlibrary{calc, quotes}
\usepackage[cache=false, outputdir=build]{minted}
\usepackage{placeins} % prnovides \FloatBarrier
\usepackage{xurl}

\definecolor{bg}{rgb}{0.9,0.9,0.9}
\newcommand\numberthis{\addtocounter{equation}{1}\tag{\theequation}}
\newcommand{\bra}{}
\newcommand{\kom}{}
\newcommand{\brak}{}
\newcommand{\braket}{}
\newcommand{\intx}{}
\newcommand{\intr}{}
\newcommand{\kla}{}
\newcommand{\klam}{}
\newcommand{\klammer}{}
\newcommand{\oper}{}
\newcommand{\tecxt}{}
\newcommand{\Bra}{}
\newcommand{\Ket}{}
\newcommand{\I}{\text{i}}
\newcommand{\e}{\text{e}}
\newcommand*{\Fourier}{\textsc{Fourier}\ }
\newcommand*{\F}{\mathcal{F}}
\newcommand*{\sinc}{\text{sinc\,}}
\newcommand{\conv}{\mathop{\scalebox{3.5}{\raisebox{-0.38ex}{\makebox[0.5\width][c]{$\ast$}}}}\limits}

\newcommand*{\titel}{Elektrostatik Kugeln}
\newcommand*{\autor}{Joachim Schwardt}

\hypersetup{pdfauthor={\autor}, pdftitle={\titel}}

%\titlehead{titlehead}
%\subject{subject}
\title{\titel}
%\author{\autor}
\date{\today}

\begin{document}

%\begin{titlepage}
\maketitle  % Titel setzen
%\tableofcontents  % Inhaltsverzeichnis setzen
%\end{titlepage}

\section{Besonders unangenehme Aufgaben}
\subsection{Oberflächenladung in einer leitenden Kugelschale}
Es ist $\rho(r)=\sigma\delta(r-a)$, und damit die Gesamtladung
\begin{align}
Q(r) = \int_0^r\rho(x) \,4\pi x^2\mathrm{d}x =\begin{cases} 4\pi\sigma a^2 & r > a \\ 0 & \mathrm{sonst}\end{cases}
\end{align}
Die Kugelschale beginnt bei $r=b$ und hat Dicke $d$. 
Das Potential sei $\phi(\infty) = 0$. \\
Für $r > b+d$:
\begin{align}
E(r) &= \frac{Q}{4\pi\epsilon_0 r^2}\\
\phi(r) &= \intx_r^\infty E(x) \mathrm{d}x = \frac{Q}{4\pi\epsilon_0 r}
\end{align}
Das Potential innerhalb der Kugelschale muss konstant sein, da das $E$-Feld verschwindet. Da das Potential stetig sein muss, gilt also für $b<r\le b+d$:
\begin{align}
E(r) &= 0\\
\phi(r) &= \frac{Q}{4\pi\epsilon_0 (b+d)}
\end{align}
Vor dem spannenden Teil noch kurz der einfache.
Für $r\le a$ gilt einfach $E=0$, da kiene Ladung mehr vorhanden ist. Das Potential ist somit konstant, und sei $\phi(r\le a):=\phi_\mathrm{inside}$.\\
Für $a<r\le b$ enthält die Gauss'sche Oberfläche jetzt die Ladung $Q$.
Dementsprechend ist 
\begin{align}
E(r) &= \frac{Q}{4\pi\epsilon_0 r^2}
\end{align}
Man könnte jetzt naiv erwarten, dass 
\begin{align}
\phi_\mathrm{naiv}(r) = -\frac{\mathrm{d}}{\mathrm{d}r}E(r)=\frac{Q}{4\pi\epsilon r}
\end{align}
ist. Das stimmt allerdings nicht ganz! Sonst wäre ja $\phi_\mathrm{naiv}(r=b)=\frac{Q}{4\pi\epsilon_0 b}$, wobei doch schon vorher $\phi(r=b)=\phi(r=b+d)=\frac{Q}{4\pi\epsilon_0 (b+d)}$ war.
Hier muss man also etwas vorsichtiger sein. Letztlich liegt der Fehler gerade in einem konstanten Offset. Das lässt sich aber eher einsehen, wenn man $\phi$ formaler bestimmt. Also für $a<r\le b$
\begin{align}
\phi(r)&=\int_r^\infty E(x)\,\mathrm{d}x=\int_r^b \frac{Q}{4\pi\epsilon_0x^2}\,\mathrm{d}x+\int_b^{b+d}0\,\mathrm{d}x+\int_{b+d}^\infty \frac{Q}{4\pi\epsilon_0 x^2}\,\mathrm{d}x=\\
&=\frac{Q}{4\pi\epsilon_0}\left( \frac{1}{r}-\frac{1}{b}+\frac{1}{b+d} \right)
\end{align}
Jetzt passt alles. Für $r=b$ ist dieses Potential wie erwartet gleich dem auf der Außenseite der Kugeloberfläche.
Für $r=a$ finden wir jetzt 
\begin{align}
\phi_\mathrm{inside}&=\phi(r=a)=\frac{Q}{4\pi\epsilon_0}\left( \frac{1}{a}-\frac{1}{b}+\frac{1}{b+d} \right),
\end{align}
also ein deutlich umständlicheres Ergebnis als das naive $\phi_\mathrm{naiv}(r=a)=\frac{Q}{4\pi\epsilon_0a}$.
\\
Hier noch der Link, der hat mir zumindest geholfen, wieder den Faden zu finden...
\url{https://physics.stackexchange.com/questions/326767/deriving-potential-inside-a-conductive-neutral-shell-containing-a-charge}

%\begin{thebibliography}{9}
%\bibitem{example} description, \url{https://www.exmapleurl}, day.\,month~2021.
%\end{thebibliography}

\end{document}